\subsection{Overview}

The literature on agricultural subsurface (tile) drainage converges on a shared meta-conclusion: despite decades of field and modeling work, the interactions among crop yield, nitrogen and phosphorus losses, and the field water balance remain insufficiently characterized across the full range of soils, climates, cropping systems, and drainage management practices that occur in tile-drained agriculture. Authors repeatedly frame future needs around \emph{extending scope} --- to more sites, soils, climates, crops, and longer observation periods --- and around \emph{sharpening the tools} used to synthesize and extrapolate, namely process-based models, monitoring networks, remote sensing, and integrated economic frameworks. \textcite{Ma2025} distil this program broadly, arguing that more research is needed to comprehensively understand the impacts of conservation practices on agricultural productivity, nutrient-loss reduction, ecosystem services, and their tradeoffs in tile-drained agricultural ecosystems. The strongest recurring lines of evidence pointing forward are: (i) the geographic and pedoclimatic narrowness of existing field trials --- concentrated in the U.S. Midwest --- which limits transferability \parencite{Zhao2016,Birru2023}; (ii) persistent structural weaknesses in the crop, cover-crop, drainage, and nutrient modules of biophysical models \parencite{Dietzel2016,Chatterjee2025,Feyereisen2006}; (iii) the need for long-term, multi-season datasets to resolve slow soil-health and climate-interaction signals \parencite{Kaspar2012,Bawa2023,Saha2024}; and (iv) the still-immature integration of biophysical, economic, and policy dimensions needed to guide adoption \parencite{Ferin2023,Malone2023,Qin2025}. This section synthesizes these open questions and research priorities rather than settled results.

\subsection{Broadening the empirical base: sites, soils, climates, and long-term observation}

A dominant theme is that existing evidence is drawn from too few sites and too short a window. \textcite{Birru2023} call for research that extends both the scope and geographic area to test and quantify the impacts of multiple cover-crop species under diverse management and biophysical conditions, and for studies spread over more sites with varied environmental conditions and management practices so that cover-crop benefits can be optimized without compromising crop-production goals. \textcite{Zhao2016} note the unequal geographic distribution of historical drainage studies and call for additional research across other states of the U.S. Midwest. Similar calls to extend testing to additional soils, climates, and management arise from \textcite{Grewal2025} (across different soil types and rainfall levels), \textcite{Negm2014} (different climatic conditions, soil types, crop rotations, farming practices, and water-management scenarios), \textcite{Saha2024} (longer periods and additional soils in the Northern Great Plains to identify systems that benefit both productivity and environmental health under a more extreme climate), \textcite{Takeshima2023} (different rainfall environments to confirm the interaction between subsurface drainage and planting pattern), and \textcite{Liu2025} (extending the experimental period and validating pipe burial depth effects on irrigation water productivity across diverse soil properties and climates, plus exploring multiple rotation systems). \textcite{BENOIT1967} likewise stress that more complete quantification of soil-water/plant-survival relationships is needed for extrapolation to other soils.

The call for \emph{long-term} experimentation is equally strong. \textcite{Kaspar2012} identify a need for additional long-term studies of the effectiveness of winter cover crops over time, and \textcite{Bawa2023} argue that a longer-term study is needed to observe the positive influence of a rye cover crop on soil-health parameters. \textcite{MartinezFeria2016} seek understanding of the long-term consequences of including rye cover crops in Midwestern systems and of the biotic and abiotic mechanisms by which rye affects yield and environmental performance. \textcite{Kumar2014b} note that long-term monitoring of fluxes under different crop rotations is needed because of seasonal climatic fluctuations. Several teams describe active or planned long-term programs: \textcite{Qi2025b} report ongoing rice--wheat and rice--oil rotation experiments to further assess subsurface-drainage performance over multiple seasons and climates; \textcite{Liu2025b} and \textcite{Liu2025c} call for long-term monitoring of subsurface pipe drainage and seepage-drainage effects across diverse environmental conditions and cropping systems; and \textcite{Thoma2005} caution that four years of data are insufficient to resolve whether higher residue cover minimizes pollutant losses and affects grain yield over time. \textcite{Acharya2019b} similarly call for a more detailed and longer-term experiment to understand climate--crop-management interactions.

\subsection{Cover crops: agronomic, environmental, and modeling frontiers}

Cover crops --- especially winter rye --- attract the largest concentration of open questions. On the agronomic side, \textcite{Waring2020} prioritize tools to promote rye growth in colder regions and study of long-term no-till in that climate, while \textcite{Bawa2023} highlight the need for management strategies to establish rye under unusually high or low temperature and precipitation and limited growing degree days. \textcite{Dougherty2020} flag that observed yield drag may reflect poor planter performance through cover-crop residue, warranting investigation. \textcite{Kaspar2007} suggest improving rye management by evaluating more cultivars, adjusting timing relative to planting, and improving establishment methods to reduce year-to-year variability and yield impacts, and note that establishment costs, logistical conflicts, and seed availability must be overcome before widespread adoption; permanent vegetation strips might likewise be improved through different species, wider strips, or altered management. \textcite{Dinnes2002} call for cultural and plant-breeding research to improve cover-crop compatibility, establishment, management, and reduced negative impacts on subsequent row crops, and \textcite{Liu2019} echo the need to improve establishment and cold hardiness. \textcite{Qi2011b} and \textcite{Qi2011} identify kura clover as requiring better agronomic management (living mulch with corn) to avoid corn yield loss under Iowa conditions.

On the environmental-mechanism side, \textcite{Lacey2015} call for research on the release of cover-crop-scavenged nitrogen in spring and on synchronizing N release from residue with the growth stages of corn and soybean. \textcite{Roth2018} identify determining the nitrogen content of grass cover crops through a producer-adoptable method as a research priority, and note that seeding rates could be adjusted while maintaining environmental benefits and that cover-crop and nutrient-management strategies could be developed to maintain cash-crop yields. \textcite{Askar2023} call for understanding the combined effects of cover crops and manure-rate application on subsurface dissolved reactive phosphorus loads and concentrations, evaluating climatic, soil, topographic, management, and crop-species factors. \textcite{Meyer2019} highlight the need to optimize trade-offs between the services and disservices of cover crops for water balance and to better understand the factors conditioning cover-crop effects on drainage. \textcite{Liu2019} more broadly call for coordinated common experiments across geophysical conditions to quantify winter crop-cover phosphorus contributions to runoff or leaching, to identify reliable plant and soil P analytes, to understand plant--root--soil--management--climate interactions on P losses, to evaluate water-quality tradeoffs under field settings, and to predict future-climate influences.

Cover-crop \emph{modeling} is a distinct frontier. \textcite{Feyereisen2006} enumerate needed advances: additional validation of the rye growth model, field verification of climate and soil-fertility effects on the ratio of rye N uptake to biomass, adjustment of radiation-use efficiency for atmospheric conditions, and better algorithms for plant N uptake. \textcite{Malone2007c} note that determining correct parameter values for winter cover crops requires thorough testing with several years of field-specific data, and that a more realistic N-uptake routine may improve corn-yield predictions. \textcite{Malone2014b} state that modeling germination and establishment of winter rye needs further research and that further analysis is the next step to assess regional cover-crop impacts. \textcite{Chatterjee2025} report that calibration of soil-organic-matter and drained-upper-limit parameters is critical and that a dedicated crop module for simulating winter rye could improve prediction of aboveground biomass and its consequences for cash-crop yield and water and N flux. \textcite{Dietzel2016} identify weaknesses in APSIM related to cover-crop effects on following-corn yields, winter soil temperatures under frozen conditions, and subsurface drainage in soybeans during flooding. \textcite{Gupta2022} propose developing a chemical-termination routine to enhance cover-crop impact studies, and \textcite{Gupta2023b} call for improving the drainage and soil-N mineralization/immobilization modules and studying extreme-event effects on cover-crop impacts to plan field-scale adaptive strategies. \textcite{Strock2004} report research in progress to predict cover-crop establishment and growth in Minnesota. \textcite{Singh2023} suggest investigating additional cover crops such as radish and legumes in hydrology and water-quality simulations.

\subsection{Drainage system design, water-table management, and materials}

Optimizing drainage design for the joint goals of crop yield and water quality is a persistent open problem. \textcite{Kladivko2021} argue that more field studies coupled with modeling are needed to develop site-specific recommendations and to optimize drainage design for both crop yield and water-quality purposes, calling continued efforts on systems approaches "urgently needed." \textcite{Zhao2016} suggest research into alternative drainage spacing aimed at balancing crop production against decreased drainage losses. \textcite{Singh2006b,Singh2007b} recommend field experiments in Iowa and the Upper Midwest to characterize how shallow and controlled drainage affect the complete water balance and crop production, including potential water-movement pathways, and note that a flexible scheme managing the control structure by soil-moisture status might reduce controlled-drainage impacts on relative yields and surface runoff. \textcite{Nash2015} state that more research is needed to determine whether managed drainage can consistently produce greater corn yield than free drainage, with benefits on claypan soils perhaps concentrated in higher-yield-potential years. \textcite{Nelson2011} call for further evaluation of drainage water management effects on high-yielding soybean cultivars, insect and disease management, cost-effectiveness in corn--soybean rotations, and grain-quality differences on claypan soils. \textcite{Dinnes2002} prioritize low-cost drainage-control and water-table-management designs that can be retrofitted to existing tile lines, and research on wetlands, buffers, and bioreactors to address residence time, carbon sources, land area, and cost limitations.

Cold-climate and peatland drainage raise specific needs: \textcite{Li2025} call for greater attention to optimizing winter subsurface-drainage and water-table management for croplands as temperatures warm, and \textcite{Läpikivi2025} identify needs to evaluate subsurface irrigation for maintaining shallow water-table depth in fields drained to current standards, to study water availability in catchments shared by multiple peatlands, and to map the distribution of cultivated peatlands and their water-table depths for catchment-wide application. \textcite{STEINHARDT1974} caution that under severe frost, subsurface drainage might irreversibly affect ploughed-soil structure and that other criteria (no-traffic systems, aeration-sensitive crops) may sometimes rule.

Drainage \emph{materials and installation} also feature. \textcite{Ghane2022} recommend long-term field investigations of knitted-sock envelope clogging across soil types to inform drainage design. \textcite{Qi2025b} report developing an integrated excavation-and-burial machine to reduce labor costs and call for investigating alternative rot-resistant straw-bundling techniques and burial methods to improve drainage longevity and waterlogging mitigation, while \textcite{Lu2018} plan to study the long-term impact of crop straw as subsurface-drainage material on drainage-water quality and soil properties under different irrigation schedules. \textcite{Ellafi2021} call for evaluating drainage designs' impacts on crop production, salinity control, and waterlogging using DRAINMOD, and for calculating predicted income across cropping systems. \textcite{Takeshima2023} recommend developing countermeasures against unpredictable rainfall, such as combining subsurface drainage with raised-bed planting.

\subsection{Nitrogen and nutrient management priorities}

Refining nitrogen management to reduce drainage losses without sacrificing yield is a core research need. \textcite{Jaynes2013} call for further field research or crop-modeling to pinpoint where between the V12 and V16 growth stages sidedress N ceases to be beneficial, and \textcite{Jaynes2015} note that sidedressing can be combined with soil tests or crop N-deficiency sensing for real-time N-rate adjustment. \textcite{Maas2022} recommend comprehensive tools integrating reliable weather forecasts, historical land management, local field and soil parameters, and crop-variety information to estimate optimum N rates, propose regional long-term agronomic-optimum-N-rate (AONR) means derived from fertilizer type, yield, and drainage status via surveys or census data, and call for study of the planting-date/agronomic-efficiency/AONR relationship and of the drivers of uncertainty and variability in agricultural systems. \textcite{Dinnes2002} prioritize better crediting of non-fertilizer N sources through soil tests, improved in-season soil-N monitoring for split applications (including earlier-season tests and models using soil respiration and meteorology), and reliable variable-rate N models incorporating remote sensing and hyperspectral approaches. \textcite{Cassman2025} propose low-cost proxy variables for monitoring soil N and water supply --- yield without applied N as a field-scale estimate of soil N-supply capacity, and partial factor productivity of applied N and rainfall as proxies for apparent N-use efficiency and economic return per unit rainfall.

At larger scales, \textcite{Zhao2016} call for constraining the magnitude and variability of crop-rotation impacts on yield-scaled N loss and for accounting for N losses on a rotation rather than annual basis. \textcite{Raymond2012} call for additional experimental and modeling work to resolve crop-management/N-export relationships, a more spatially intense monitoring system of riverine outflow and net anthropogenic nitrogen inputs coupled with detailed crop-management (particularly tile-drainage) information, and exploration of alternative drainage approaches, cover crops, and crops with more active rooting and transpiration. \textcite{Pasley2021} note very little research on more desirable rotation crops in simple Midwestern two-crop rotations from a leaching perspective. \textcite{Deng2011} plan to apply revised DNDC to more watersheds, improve parameterization for N-loss prediction, test the variance of surface-cover and management factors across crop-growth stages, and schedule aquatic biogeochemistry to address the fate of N in streams, rivers, lakes, reservoirs, and estuaries. \textcite{Kanwar1997} reinforce the need for studies on chemical placement, rate, and timing for various tillage practices to reduce tile-drainage chemical losses, incorporated into at least a two-year corn--soybean rotation to ensure economic and environmental viability. \textcite{Bakhsh2004} note that overprediction of soybean yields may require refining N-cycling algorithms for N\textsubscript{2} fixation and plant N-uptake. \textcite{Malone2018} call for field studies to evaluate fertilized/harvested rye cover crop, which if confirmed could offer a strategy to raise producer revenue and net energy production while reducing drainage N loss.

\subsection{Irrigation, salinity, and water reuse in arid and saline systems}

A parallel body of work addresses drainage in irrigated, arid, and saline settings. \textcite{Yao2017} call for long-term geohydrological and crop data --- including intensive electromagnetic-induction surveys, groundwater-dynamics time series, and crop biomass --- to validate simulations across irrigation/drainage systems and agronomic patterns. \textcite{Alzraiee2013} argue that future modeling must systematically evaluate alternative field-scale irrigation and drainage strategies to manage waterlogging and high soil salinity under ambiguity, that a statistical notion of the likelihood of achieving environmental and crop-production goals field-by-field would guide riskier watershed-scale plans, and that considering multiple irrigation seasons might reveal different spatial distributions of yield impact. \textcite{Bethune2008} recommend that, where tubewell drainage manages deep percolation, land-use planning steer irrigation development toward low-salinity, high-yielding aquifers, encourage sprinkler irrigation, and move new development to areas where salinity-management costs are low. \textcite{Qadir2001} suggest incorporating short-duration and salt-tolerant crops (e.g., Sesbania, Bermuda or Kallar grass) into reclamation rotations, and \textcite{Castanheira2010} call for quantifying rainfall amounts and irrigation-efficiency gains with respect to salt-removal capacity and reduced salt load.

Water-productivity and technology questions also recur. \textcite{Feng2019} propose using their outputs to design irrigation schedules under different hydrological years and soil water-salt environments and using the coupled HYDRUS--EPIC model to evaluate the future sustainability of saline-water irrigation. \textcite{Jia2017} recommend measuring water-table changes and their contribution to storage and evapotranspiration, recomputing water-application efficiency per day between soil samplings to account for water-table rise, reformulating the efficiency equation, and using longer irrigation durations or slower flow rates to improve uniformity. \textcite{Qin2025} call for scalable data-management infrastructures for long-term field deployment (energy-efficient edge computing, low-power IoT protocols), life-cycle cost and return-on-investment assessment under varying market and climate conditions, integration of agronomic models with environmental and economic frameworks in a unified precision-agriculture platform, and low-temperature crystallization to treat brine toward near-zero discharge. \textcite{Abdelkrim2021} note that specific irrigation dates and amounts should be scheduled using criteria beyond modeling, such as farmers' expert knowledge and sensors. \textcite{DINAR1989} call for a more complete economic evaluation of shifting irrigation technology including associated management variables, and \textcite{Qian2025} highlight the need to study solute transport during cultivation from a components perspective and to test applicability under drip irrigation. \textcite{ElHassanin2022} suggest that established thresholds could seed large-scale modeling or the development of water/drainage-management measures to raise field water-use fraction, crop-water productivity, and wheat production, and \textcite{RADSTAKE1988} recommend continued monitoring of crop yield and water-quality change to improve forecasts of groundwater suitability for irrigation in the Nile Valley. \textcite{Ayars2006} note that extending their method to other crops will require additional studies, and \textcite{Karegoudar2022} advise that policymakers may proceed under all three water-availability scenarios. Forage-water-use comparisons remain open: \textcite{Mundy2006} judge a comparison of the water-use efficiency of forages under flood versus spray irrigation worthwhile, and \textcite{Greenwood2006} defer forage yield and plant--soil--water responses to companion papers.

\subsection{Advancing process-based and biophysical models}

Beyond cover-crop and drainage modules, authors call broadly for stronger, better-validated, and more transferable models. \textcite{SKAGGS1994} state that much remains to be learned about the mechanisms governing pollutant losses from drained soils, that research must improve process understanding before quantification is possible, that methods are needed to quantify how alternative designs and operations affect crop yields and environmental impacts while considering interdependent cultural, fertility, and pest-management practices, and that better datasets are needed to test models across wide ranges of soil, crop, and climate. \textcite{Negm2014} call for more intense field measurements to validate their model and additional testing across conditions, noting its potential for simulating long-term climate-change impacts. \textcite{Grewal2025} prioritize improving prediction of high-load peak events, refining the concentration component of the model, further testing the fertilizer-effect hypothesis under ponding, and empirical data to validate predictions for cover-crop and controlled-drainage practices. \textcite{Gunn2018} report that DRAINMOD is being coupled with a crop model (DSSAT) to capture other factors influencing maize relative yield under subirrigation. \textcite{Tian2012} call for field verification of soil-physical-property changes, investigation of early-successional plant growth and its hydrological effects, verification of temporal stomatal-conductance changes against field measurements, and cautious evaluation of adding understory-growth and prescribed-burning modules. \textcite{Morrison2014} seek greater sensitivity to low flows, yearly crop-information inputs for calibration/validation consistency, and a macropore-flow algorithm to capture poorly draining soils. \textcite{Smith2008} argue that a validated process-based rather than empirical model should be applied to capture complex variable interactions. \textcite{Ndulue2024} recommend incorporating CO\textsubscript{2} effects, diseases, nutrient availability, cultivar type, and planting date for greater confidence in climate-change results. \textcite{Ma2025} argue that systematic data collection on water fluxes, solute-nutrient fluxes, greenhouse-gas emissions, root development, and aboveground biomass would facilitate model validation and refine mechanisms and parameterization. \textcite{Li2025} note that winter-crop-growth simulations need additional crop-genetic-parameter calibration and must account for elevated CO\textsubscript{2} and alternative management, and that simpler, faster biophysical models may be viable for large-scale winter-cropland simulation. \textcite{Dietzel2016} additionally call for improved subsurface-drainage simulation within soybean in APSIM. \textcite{Qi2010b} report a field-scale experiment being conducted in parallel to a non-weighing lysimeter study to validate rye cover-crop effects on soil-water dynamics and nitrate leaching, and \textcite{Malone2007c} note that per-plot soil-parameter calibration would improve drainage predictions.

\subsection{Climate change, extremes, and system resilience}

Climate adaptation frames many priorities. \textcite{Saha2024} stress continuing to identify practices and combinations that increase the resilience of agricultural systems to climate extremes. \textcite{Malone2020} lay out a modeling agenda addressing additional simulation periods, other representative concentration pathways, management options (reduced fertilizer rates; earlier cash- and cover-crop harvest, termination, and planting), cash-crop cultivars, cover-crop types, and additional soils and locations. \textcite{Ferin2023} recommend approaches allowing within-county variation in fertilizer application rate, crop-specific manure applications as data become available, multiple global-climate-model forcings from CMIP5 or CMIP6, and attention to co-benefits of perennial grasses such as carbon sequestration and lower N\textsubscript{2}O loss. \textcite{Dietzel2016} argue that efforts to raise water-use efficiency should design cropping systems that function under both drought and flood, diversify whole-farm designs to include drought- and flood-tolerant plants, and pursue engineering solutions such as capturing excess floodwater for later use in drought. \textcite{PeltonenSainio2021} call for more information on the costs and benefits of adaptation measures for changing precipitation, a thorough assessment of the payback period for irrigation investment in Finland, stronger knowledge-sharing on carbon farming and soil-carbon maintenance, and additional funding and decision tools for both irrigation and drainage investment. \textcite{Morton2015} identify a need for increased localized understanding of climate patterns and how farmers incorporate actual and perceived experiences of excess precipitation into decisions, more accurate downscaled climate information and local forecasts beyond 3--5 days, and better understanding of the relationships among practices and adoption timescales. \textcite{Gupta2023b} likewise call for studying extreme-event effects on cover-crop impacts to plan field-scale adaptation.

\subsection{Remote sensing, precision tools, and decision support}

Emerging sensing and data technologies are widely seen as enabling the next phase of research. \textcite{Khanal2020} outline a broad agenda: determining the optimal field conditions and image quality for detecting tile-drainage patterns and advanced processing for automatic extraction of tile-drain lines; downscaling techniques combining multiple data sources for finer-resolution soil moisture; integrating remote sensing with machine learning to assess soil properties and compaction; coupling crop growth stages with forecasted weather for supply forecasting; leveraging high-resolution remote sensing with crop modeling to understand within-field nitrogen-stress variability and variable-rate N benefits; mapping crop genotypes and their role in grain quality; developing databases for early crop-disease detection; applying semi-supervised learning and generative adversarial networks to weed identification; and site-specific residue management. \textcite{Cassman2025} propose a robust spatial framework combining climate and soil type with geo-referenced farmer-reported rotation, yield, and management data, tailoring a Climate Zone and Technology Extrapolation Domain framework to delineate conditions under which soil-health-improving-practice adoption yields similar impacts, and using remote-sensing, satellite-imagery, and on-farm databases to estimate field-specific benefits. \textcite{Reid2018} call for validating individual components and aggregate edge-of-field losses for Canadian conditions, a large effort to quantify edge-of-field-to-surface-water delivery ratios, improved input-data collection with possible remote sensing of topography, soil moisture, crop area, and soil cover, and intensity-based indicators of phosphorus loss per unit production developed with commodity organizations. \textcite{Julien2023} envision a soil-wetness-index approach forming part of a decision-support tool integrated into marketing-authorization recommendations, while noting that assumptions about drainage-season insensitivity to winter crop type and site-dependent thresholds need testing across soils.

\subsection{Cropping-system diversification: perennials, bioenergy, and rotations}

Diversified and perennial systems raise distinctive research needs. \textcite{Arrueta2022} call for establishing a clear causal link between perennial-crop implementation and nutrient loads across fields with varying physical and chemical soil characteristics, and note that government policies encouraging landscape diversification and rewarding environmental services may be needed to overcome adoption barriers. \textcite{Wilson2023} call for robust observed water-quality datasets for perennialized cropping systems to understand their regional water-quality impact, supporting watershed-scale learning and collective action around novel crops. \textcite{Guo2018} propose large-scale assessment of bioenergy-crop scenarios in the Mississippi River system and other regions (e.g., China's marginal lands), quantification of biofeedstock production and its water-quantity/quality impacts in mildly sloped watersheds, investigation of pasture loss during bioenergy planting, field experiments to quantify water-quality impacts, and monitoring of nutrient-reduction efficiencies with cost/benefit assessment to alleviate Gulf of Mexico hypoxia. \textcite{Gentry2025} report future work refining management within a diverse cropping system, documenting economic viability against the conventional two-year rotation, nested microplot studies of wheat varieties and seeding rates for double-crop soybean, and planned water-infiltration studies of earthworm influence on tile flow and nitrate loss. \textcite{Dinnes2002} call for perennial cash-cropping systems that maintain economic returns while reducing nitrate leaching. \textcite{Malone2018} highlight double-cropping fertilized rye as a promising revenue, net-energy, and drainage-N strategy pending field confirmation, and \textcite{HUNDAL1976} report ongoing studies of meadow cropping effects on soil conditions and grain yields.

\subsection{Economics, policy, adoption, and the human dimension}

Bridging biophysical findings to on-farm and policy decisions is a recurrent frontier. \textcite{Waring2020} prioritize the economics of cover crops. \textcite{Roth2018} note that quantifying soil-health impacts, surface and subsurface phosphorus losses, and denitrification changes could increase cost recovery, and that a long-term average of total cover-crop cost recovery could guide policymakers on subsidies. \textcite{Mourtzinis2021} observe that a long-term cumulative yield-and-profit assessment of drainage installation would require installation cost, yield responses in other rotation crops, land value, and other farm benefits, and could quantify drainage adoption's contribution to soybean yield gain. \textcite{Malone2023} call for information on producer participation, pilot programs in key locations, field trials and modeling clarifying how primary-crop yield, rye yield, and N losses vary with soil, climate, and management, and additional studies to advance policies optimizing environmental quality, bioenergy, carbon, yields, and profits. \textcite{Ferin2023} argue that future efforts should ultimately include the multitude of biophysical factors and socioeconomic nuances through integrated biophysical--economic modeling at the field scale. \textcite{Kröger2012} urge committing significant resources to replicable studies in the Lower Mississippi Alluvial Valley to estimate best-management-practice effectiveness under realistic conditions, plus additional field-based monitoring to generate efficiency estimates for conservation-program and policy analysis. \textcite{Bakhsh2007b} argue that landscape/hydrologic spatial zones should be managed with suitable tillage and cropping systems for watershed sustainability and that decision tools should account for soil and landscape effects on yields. \textcite{Nelson2011} include cost-effectiveness in corn--soybean rotations among evaluation needs, and \textcite{Jaynes2001} conclude that dramatic changes in management on artificially drained soils will be required before corn production can be considered sustainable.

\subsection{Soil health, biology, and physical processes}

Slower-acting soil and biological processes are flagged as understudied. \textcite{Han2017} note that the link between \emph{Fusarium} abundance and soybean root rot needs examination and that long-term drainage may affect field \emph{Fusarium} populations and, under certain conditions, soybean-root health and growth. \textcite{Kumar2014b} call for measuring biogeochemical processes involved in methane production and oxidation and modeling these fluxes from poorly drained soils. \textcite{Shipitalo2004} suggest that disrupting the linkage of cracks to earthworm burrows through the plow pan --- via conservation tillage as an alternative to moldboard plowing and lightweight tractors to reduce compaction and prevent plow-pan reformation --- might reduce sediment and contaminant movement to subdrains, noting no-till has reduced sediment losses in crack-prone clayey soils. \textcite{Daigh2014} call for precisely quantifying evapotranspiration to explain drain-flow differences among cropping systems, and \textcite{Tian2012} note early-successional growth and soil-physical-property changes need verification. \textcite{DarziNaftchali2017b} argue that the drought tolerance and phenological traits of rice cultivars should guide selection of drainage systems and drying periods, timed to cultivar sensitivity to water deficit. \textcite{Nash2017} advise farmers to test forage sorghum for toxic nitrate levels before grazing regardless of perceived drought. \textcite{Manninen2018} call for understanding the reactions and role of agriculturally derived dissolved organic matter in marine systems over time, and \textcite{GiménezValero2024} anticipate complementing metabolomic fruit analysis with studies of cultivation-technique effects on citrus fruit quality and size.

\subsection{Watershed-scale synthesis, monitoring networks, and data infrastructure}

Finally, several authors emphasize scaling from field to watershed and building the data infrastructure to support it. \textcite{Baffaut2020} propose modifying the Soil Vulnerability Index to differentiate surface from subsurface artificial drainage (adjusting the risk class of the runoff or leaching component), to consider the leaching component both with and without artificial drainage to distinguish permeable-soil from drainage causes, to account for irrigation impacts on surface-water quality and drainage, to incorporate rainfall intensity (possibly via the R-factor), and to determine exact soil components where complex soils or land leveling complicate classification. \textcite{Singh2023} suggest calibrating and validating models for multiple sub-watersheds and comparing modeled versus observed crop yield at the sub-watershed level for more reliable estimates. \textcite{Smith2008} recommend focusing future conservation-assessment efforts on managing potholes in cropped fields --- especially corn fields where potholes drain directly to ditches --- and their effects on tile-drainage water. \textcite{Guo2018} extend this to large-scale bioenergy assessment across river systems. \textcite{Ferin2023} and \textcite{Deng2011} both point toward multi-forcing, multi-watershed application. Underlying all of these is the data imperative articulated by \textcite{Ma2025} and \textcite{SKAGGS1994}: better, longer, and more spatially representative datasets --- spanning water and solute fluxes, greenhouse-gas emissions, root and biomass development, and management history --- are the foundational requirement for validating models, resolving tradeoffs, and delivering the site-specific, climate-resilient drainage-management guidance the field ultimately seeks.